\section{Additional Robustness}
\label{app:robustness}

\subsection{Wild cluster bootstrap}
With 449 date clusters, cluster-robust inference on the headline
reduced form is unlikely to be distorted, but the analysis plan locks a
wild cluster bootstrap as a check. Using Webb weights and 9{,}999
replications on the \emph{corrected} weekday step in the \$2.50
regime---the preferred specification, not the uncorrected one---the
bootstrap rejects the null of no effect at $p < 10^{-6}$, with a 95
percent confidence set of $[0.2834, 0.3012]$ around a point estimate of
\$0.2923 \parencite{cameron2008bootstrap}.

\subsection{Permutation over candidate cutoffs}
The clock design has a natural randomization test: re-estimate the
discontinuity at candidate minutes where no surcharge changes and locate
the true cutoff in the resulting distribution. Every candidate---true
and placebo alike---is estimated with the same specification and the
same $\pm 30$ minute bandwidth, so the placebo distribution is
comparable to the actual statistic; candidates are drawn only from
windows that never cross 4:00pm, so no placebo inherits part of the true
discontinuity. Of 60 such candidates, \emph{none} produces an absolute
estimate as large as the true cutoff's $0.2401$.

\subsection{Placebo cutoffs}
Placebo cutoffs at 3:15pm and 5:00pm---one before the surcharge begins,
one after it is already in force---produce first stages of $-0.001$ and
$-0.023$ respectively, against $1.951$ at the true cutoff, and tip-side
effects of $-\$0.0003$ (SE \$0.0068) and $-\$0.0079$ (SE \$0.0071).
Both sides of each placebo are null. Whatever generates the 4:00pm
discontinuity is specific to the minute at which the fare schedule
changes.

\subsection{Mechanical benchmark}
A purely mechanical account of pass-through predicts a tip step equal to
the surcharge step times the prevailing tip rate on the base. In the
\$2.50 regime the surcharge step is \$1.951 and the pre-cutoff tip rate
is $0.174$, for a mechanical prediction of \$0.339. The corrected
estimate is \$0.292---about 86 percent of the mechanical benchmark. The
remaining fourteen percent is the behavioral offset: the small declines
in exact-default adherence and the small rise in zero-tipping documented
in Section~\ref{sec:results}. Pass-through is therefore close to, but
slightly below, complete.

\subsection{Synthetic difference-in-differences}
Because the binary toll event study fails a strict pre-trend test, we
re-estimate the adoption effect with synthetic difference-in-differences
\parencite{arkhangelsky2021synthetic}, which reweights both control
zones and pre-periods to match the treated group's pre-treatment
trajectory---the estimator built for exactly this failure mode. On a
balanced zone-month panel of fare-conditioned tips (treated: the 52
CBD-pickup zones; donors: never-eligible zone pools), the SDID average
treatment effect is \$0.061 (bootstrap SE \$0.050), an implied
$\rho = 0.102$---inside the family of estimates but underpowered at
zone-level aggregation, where the design surrenders the cell-level
precision of the main specifications. An in-time placebo assigning a
fake July 2024 adoption returns $-\$0.125$ (SE $0.069$), reflecting the
same 2024 drift the Wald test flags. We read SDID as corroborating the
magnitude while confirming, rather than repairing, the pre-period's
imperfection; the dose staircase remains the spatial design's strongest
evidence.
